\section{Setup}
\label{sec:setup}

\paragraph{Task and agents.} A Llama-3.2-1B therapist (bf16, LoRA) conducts Motivational
Interviewing against a simulated client. The client and the reward oracle are both
\primaryjudge{}; the client is conditioned on one of 96 distinct persona prompts varying the
presenting problem and the degree of cooperativeness. The training reward is the mean of two
published LLM-evaluator rubrics for session satisfaction and working alliance, scored on the
conversation so far.

\paragraph{Optimiser.} Both arms use the same iterative preference-tree procedure. At each
therapist turn the current policy samples $M{=}8$ candidate completions; each is scored by the
oracle; the best and worst form a preference pair when their scores differ by more than a fixed
$\tau$; the pairs train a DPO update; the resulting adapter generates the next iteration's data.
The 96 conversations generated at the start of iteration $n$ are the evaluation set for the policy
that produced them, so evaluation needs no additional generation.

\paragraph{The manipulation.} The two arms differ in one thing. Under $K{=}0$ the oracle scores
the conversation prefix plus the candidate turn. Under $K{=}5$ the candidate is first extended by
five further simulated turns --- client, policy, client, \ldots --- using the same policy, and the
oracle scores that continuation. Everything else is matched by construction and verified against
both runs' stored run metadata: sampling temperature $1.2$, $M{=}8$, minimum context length $12$,
target conversation length, $\tau$, DPO $\beta$, epochs per iteration, and the oracle itself.

Look-ahead is therefore purely a statement about \emph{what context the oracle grades}. It changes
no loss, no hyperparameter, and no sampling distribution.

\paragraph{Budget.} Matching iterations does not match cost, and the mismatch runs against the
intervention rather than for it. Across the ten training iterations the $K{=}5$ arm built $7{,}548$
preference groups and trained on $6{,}416$, against the $K{=}0$ arm's $6{,}240$ and $4{,}935$ ---
and each $K{=}5$ candidate additionally pays five simulated turns. Iteration-matching is the
correct scientific control, but it is generous to $K{=}0$ on budget, and it rules out the reading
that the $K{=}5$ arm hacks less because it saw less training signal: it saw more.

\paragraph{Evaluation.} Every conversation of every model state is scored on eight instruments by
the training oracle, and independently by \heldoutjudge{}, a model from a different family that
never played the client and never contributed a training reward. This paper reads two of them: an
MI-inconsistency coder that counts six specific violations per session (over-praise,
advice-without-permission, directing, confronting, warning, judging) and rates their severity, and
the reward rubric itself. Because each iteration reshuffles the same 96 personas, every comparison
below is \emph{paired on the recovered persona}, not on file order; means survive a file-order
join but effect sizes and intervals do not.

\paragraph{Statistics.} All contrasts are persona-paired Wilcoxon signed-rank with Cohen's \dz{}
and Holm correction \emph{within} a family, where a family is one (arm-pair, iteration, channel
group) --- for instance the six per-session violation counts at iteration~8. Corrections are not
pooled across iterations; each matched iteration is its own family. We report per-session counts
as the primary quantity and per-turn rates as secondary, because the two arms differ in turn count
and a rate can move with the count fixed.

\paragraph{Horizon.} Both arms ran ten training iterations, so all matched comparisons run over
eleven points, iterations 0--10, with 96 persona-matched conversations per arm at every point.
